\section*{Chapitre \ref{ch:ll:protocols-i-fix} --- Protocoles I~: FIX}

\begin{solution}{exo:ll:protocols-i-fix:1}
Le corps est \texttt{35=0|} (5 octets), \texttt{49=A|} (5), \texttt{56=B|} (5), \texttt{34=2|} (5) et le champ d'heure d'envoi (25)~:
45 octets, soit \texttt{9=45}. Les octets précédant \texttt{10=} totalisent un nombre dont le reste modulo 256 vaut 75~: \texttt{10=075}, comme
\texttt{firm\_fixengine.encode} le confirme pour le numéro de séquence 2 à 14:30:00.
\end{solution}

\begin{solution}{exo:ll:protocols-i-fix:2}
Une demande de test (test request) part après 1.2 intervalles de silence, à 10:00:36~; si aucun \omterm{def:ll:protocols-i-fix:session}{heartbeat} portant son identifiant ne revient dans un intervalle de plus,
la session déclare la connexion morte à 10:01:06, soit 66 secondes après le dernier message.
\end{solution}

\begin{solution}{exo:ll:protocols-i-fix:3}
Un ResendRequest avec BeginSeqNo 52 et EndSeqNo 0~; il retient le message 57. Le courtier renvoie 52--56 (et tout ce qui suit, car
EndSeqNo 0 l'exige)~: les messages applicatifs avec PossDupFlag et OrigSendingTime, les messages administratifs remplacés par
SequenceReset-GapFill~; puis le 57 est délivré depuis la file d'attente, et la copie renvoyée du 57 est ignorée comme doublon.
\end{solution}

\begin{solution}{exo:ll:protocols-i-fix:4}
En échangeant les identifiants d'expéditeur et de destinataire du \omterm{def:ll:protocols-i-fix:session}{heartbeat} de l'exercice~1 donne \texttt{49=B|56=A|} au lieu de
\texttt{49=A|56=B|}~: les mêmes octets dans un autre ordre, donc même longueur et même somme de contrôle, 075. La somme de contrôle doit seulement
détecter les erreurs de cadrage et les altérations d'un seul octet~; les sommes de contrôle propres à TCP protègent les octets en transit, et la couche session rejette
un message adressé à un mauvais destinataire d'après ses identifiants.
\end{solution}

\begin{solution}{exo:ll:protocols-i-fix:5}
Ils décrivent la liaison au moment de leur envoi~: un \omterm{def:ll:protocols-i-fix:session}{heartbeat} dit «~je suis vivant maintenant~», une demande de test réclame une réponse
maintenant, une \omterm{def:ll:protocols-i-fix:session}{demande de renvoi} (resend request) réclame des messages qui ont pu être envoyés depuis. Renvoyés plus tard, ils seraient dénués de sens ou nuisibles~: une
demande de test périmée exigerait une réponse, une \omterm{def:ll:protocols-i-fix:session}{demande de renvoi} périmée déclencherait un second renvoi. Le \omterm{def:ll:protocols-i-fix:session}{comblement de trou} (gap fill) fait avancer la séquence
au-delà d'eux sans les rejouer.
\end{solution}

\begin{solution}{exo:ll:protocols-i-fix:6}
Sur la mesure retenue~: $200\,000 \times \qty{84}{\nano\second} \approx 1.7~\%$ d'un cœur avec la vue zero-copy,
$200\,000 \times \qty{892}{\nano\second} \approx 18~\%$ avec l'analyseur à base de map, et $200\,000 \times
\qty{5.3}{\micro\second} \approx 107~\%$ avec la référence Python, qui ne tient pas la cadence du tout.
\end{solution}

\begin{solution}{exo:ll:protocols-i-fix:7}
Dans \texttt{on\_bytes}, intercepter l'erreur de décodage de type \texttt{tag} avant le contrôle de séquence, lire MsgSeqNum dans les octets bruts
s'il est présent, et émettre \texttt{35=3} avec \texttt{45=<that number>} et un texte~; faire ensuite avancer le numéro attendu (le
message a bien été reçu, seulement pas traité). Faire de même dans le \texttt{on\_bytes} en C++ sur \texttt{Error::tag}, ajouter un message malformé
à \texttt{make\_fix\_fixtures.py}, régénérer la trace attendue et lancer les deux tests.
\end{solution}

\begin{solution}{exo:ll:protocols-i-fix:8}
Les messages envoyés pendant la coupure de la connexion (ou en transit au moment de la rupture) ne sont jamais livrés~: un rapport d'exécution pour une exécution
est perdu, et la position de la société de trading est fausse. Réinitialiser n'est acceptable qu'à condition d'avoir un autre moyen de reconstituer l'état de chaque ordre,
comme un drop copy ou des requêtes sur l'état des ordres après chaque reconnexion, ce qu'exigent les plateformes qui ne prennent en charge aucune \omterm{def:ll:protocols-i-fix:session}{demande de renvoi}.
\end{solution}

\begin{solution}{pb:ll:protocols-i-fix:1}
\begin{enumerate}
\item Environ 66 secondes~: une demande de test après 36 secondes de silence, puis un intervalle supplémentaire de 30 secondes pour sa réponse.
\item 13\,227 rapports d'exécution avec la graine fixée, et aucun \omterm{def:ll:protocols-i-fix:session}{heartbeat}~: à 200 messages par seconde, le courtier n'est jamais inactif pendant 30
secondes.
\item $13\,227 \times 205 = 2\,711\,535$ octets, soit environ \qty{21.7}{\milli\second} à 1~Gb/s.
\item Deux allers-retours (logon, \omterm{def:ll:protocols-i-fix:session}{demande de renvoi}), les octets et l'analyse~: environ \qty{24.8}{\milli\second} avec le moteur C++
et environ \qty{94}{\milli\second} avec la référence Python.
\item 5\,776 messages, dont 38 messages applicatifs et les autres \omterm{def:ll:protocols-i-fix:session}{heartbeats}.
\item Avec \omterm{def:ll:protocols-i-fix:session}{comblement de trou}~: 77 messages, soit les 38 messages applicatifs et 39 \omterm{def:ll:protocols-i-fix:session}{comblements de trou}, un par suite de \omterm{def:ll:protocols-i-fix:session}{heartbeats}. Sans, 5\,776.
\item Environ \qty{2.5}{\milli\second} avec \omterm{def:ll:protocols-i-fix:session}{comblement de trou}, contre \qty{42}{\milli\second} sans.
\item Il n'y a aucun message administratif à remplacer.
\item Ils se rapportent à la liaison telle qu'elle était au moment de leur envoi (exercice~5).
\item Vérifier s'il a déjà été traité, à partir de son identifiant d'exécution, et ne l'enregistrer qu'une seule fois.
\item À partir du drop copy de la plateforme ou en interrogeant l'état de chaque ordre ouvert après la reconnexion, et réconcilier avec les
registres propres de la société de trading avant de négocier de nouveau.
\item Ses prochains numéros de séquence sortants et entrants attendus, ainsi que les messages envoyés (pour répondre aux \omterm{def:ll:protocols-i-fix:session}{demandes de renvoi}), écrits
avant ou au moment de leur envoi.
\item \textbf{Résultat nommé.} Une coupure de 66 secondes en période chargée, à 200 messages par seconde, laisse 13\,227 messages à renvoyer, soit environ
\qty{25}{\milli\second} pour se resynchroniser avec le moteur C++ et \qty{94}{\milli\second} avec la référence Python, avec ou
sans \omterm{def:ll:protocols-i-fix:session}{comblement de trou}~; un week-end calme de 48 heures laisse 5\,776 messages, que le \omterm{def:ll:protocols-i-fix:session}{comblement de trou} ramène à 77, de \qty{42}{\milli\second} à
\qty{2.5}{\milli\second}.
\item Oui~: une détection en 2.2 secondes laisse environ 450 messages au lieu de 13\,227~; le coût est un \omterm{def:ll:protocols-i-fix:session}{heartbeat} par seconde dans chaque sens
au repos, et de fausses alertes lorsqu'un pair ralentit brièvement.
\item Les messages plus anciens que le stockage ne peuvent pas être renvoyés~: le courtier doit les combler par gap fill ou réinitialiser la séquence au-delà d'eux, et les rapports d'exécution
perdus doivent être récupérés par réconciliation.
\item Il faut mettre à jour immédiatement les positions et les états d'ordres à partir d'eux, mais les juger comme des événements passés~: les cours et le carnet d'ordres ont
évolué depuis leur heure d'envoi initiale (122).
\item Avec des conversations de référence (messages perdus, doublons, \omterm{def:ll:protocols-i-fix:session}{demandes de renvoi} dans les deux sens, réinitialisations), comme le fait la construction,
et contre le simulateur de bourse du livre~10.
\item Oui~: 200 messages par seconde à \qty{5.3}{\micro\second} représentent moins de 0.1~\% d'un cœur~; c'est la rafale de renvoi qui
prend des dizaines de millisecondes à absorber.
\item Le cas chargé~: son renvoi est le plus volumineux, et il survient lorsque les marchés bougent.
\item Garder identiques les vues de deux sociétés sur leur conversation, message par message, quoi que fasse le réseau.
\end{enumerate}
\end{solution}

\begin{solution}{iq:ll:protocols-i-fix:1}
Champs \texttt{tag=value} séparés par SOH~: BeginString, BodyLength et MsgType en premier, le reste de l'en-tête, le corps, et
CheckSum en dernier. Le récepteur lit les deux premiers champs~; BodyLength donne la longueur jusqu'au champ de somme de contrôle, qui fait toujours
sept octets (\texttt{10=nnn} et SOH).
\iqlookfor{longueur du corps pour le cadrage, somme de contrôle comme marqueur de fin.}
\end{solution}

\begin{solution}{iq:ll:protocols-i-fix:2}
Le récepteur voit un numéro de séquence supérieur à celui attendu, envoie un ResendRequest à partir du numéro attendu (EndSeqNo 0), et
retient les messages suivants. L'émetteur renvoie les messages applicatifs avec PossDupFlag et OrigSendingTime, et comble
les messages administratifs par gap fill avec SequenceReset. Le récepteur les traite dans l'ordre, libère les messages retenus, et ignore
les doublons.
\iqlookfor{mise en attente, renvoi, gap fill, doublons.}
\end{solution}

\begin{solution}{iq:ll:protocols-i-fix:3}
Repérer chaque délimiteur (vectorisé), analyser chaque étiquette sur place, et consigner (tag, décalage, longueur) dans un tableau de taille fixe~; les valeurs sont
des vues dans le tampon de réception, valables tant qu'il existe. Valider les champs de cadrage, la longueur du corps et la somme de contrôle en une seule passe.
Renvoyer la vue, ou un type d'erreur.
\iqlookfor{vues, tableau de taille fixe, et validation en une seule passe.}
\end{solution}

\begin{solution}{iq:ll:protocols-i-fix:4}
Pour distinguer le silence d'une connexion morte~: un côté inactif pendant l'intervalle envoie un \omterm{def:ll:protocols-i-fix:session}{heartbeat}~; un côté qui n'entend rien envoie une
demande de test, et abandonne s'il ne reçoit aucune réponse. Un intervalle plus court détecte les pannes plus tôt, donc moins de messages partent dans le vide, au
prix d'un trafic accru et de fausses alertes lorsqu'un pair marque une pause.
\iqlookfor{délai de détection face au trafic et aux fausses alertes.}
\end{solution}

\begin{solution}{iq:ll:protocols-i-fix:5}
Les champs binaires à décalage fixe n'exigent ni recherche, ni analyse de nombres, et consomment moins de bande passante, et des tailles fixes simplifient le cadrage~: des dizaines de
nanosecondes par message, de façon prévisible. Ils renoncent à la lisibilité, au dictionnaire commun partagé entre les plateformes, et à la souplesse~:
chaque plateforme a ses propres formats et versions.
\iqlookfor{coût d'analyse et déterminisme face à l'interopérabilité.}
\end{solution}

\begin{solution}{iq:ll:protocols-i-fix:6}
Le moteur a perdu son numéro de séquence sortant (non persisté, ou restauré depuis une vieille copie) et a redémarré sous le numéro
attendu par le courtier. Restaurer les numéros persistés et le stockage~; s'ils sont perdus, convenir hors bande d'une réinitialisation de séquence avec le courtier,
puis réconcilier l'état de chaque ordre à partir des registres du courtier ou d'un drop copy avant de négocier, en enregistrant les exécutions par
identifiant d'exécution, de sorte qu'aucune ne soit comptée deux fois.
\iqlookfor{état persisté, réinitialisation maîtrisée, et réconciliation.}
\end{solution}
